\section{Interes continuo}

\section{Inflacion}
\paragraph{Definicion} La inflacion es la perdida de poder adquisitivo de una moneda, puesto que, 
esta misma afecta a todas las personas en todas partes del mundo a todo momento es un elemento a tener en cuenta cuando hablamos de planeación financiera.\\
Es ahora que usted ya que pudo apreciar los elementos mas basicos de las finanzas por separado, es capaz de 
agregar este componente común de la realidad al analisis que realizamos.\\

Por tanto la pregunta que surge es como este fenomeno afecta a nuestros planes, una forma común de ejemplificarlo es con la siguiente situacion, 
el primero del mes un pedazo de pan vale Bs1 por cada hogaza, pero para la misma fecha el siguiente año la misma hogaza vale Bs2, 
para nosotros esto significaria que ahora para comprar la misma cantidad de pan necesitariamos el doble de dinero que antes.\\

Es asi que para cuantificar este efecto se utiliza la siguiente formula, usualmente conocida como teorema de Fisher:
\begin{center}
    $\Large{\boxed{(1+r)=\frac{(1+i)}{(1+\pi)}}}$ 
    \begin{align*}
        r \rightarrow   &\text{Tasa de rendimiento real}\\
        i \rightarrow   &\text{Tasa de interes}\\
        \pi \rightarrow &\text{Tasa de Inflacion}\\
    \end{align*}
\end{center}
Queda espacio para aclarar, para usar la formula de fisher todas las tasas tienen que estar en un mismo 
periodo efectivo, usted ya se encuentra familiarizado con este concpeto, puesto que, en capitulos anteriores 
usted ya se encontro con las tasas equivalentes, el llevar todas las tasas de la formula de Fisher a un mismo 
periodo, entonces no es mas que una simple aplicacion de estas mismas.\\
% PEDIR AUTORIZACION AL CESAR PARA PODER AGREGAR ESTA FORMULA MAS

% Aunque la formula de Fisher sea muy util para algunos movimientos simples, seria bueno que conozca esta segunda formula, la cual tiene alguna aplicacion extra que le puede ser util:
% \begin{center}
%     $\Large{\boxed{\frac{C_0}{P}(1+r)^t=\frac{C_t}{P(1+\pi)^t}}}$ 
%     \begin{align*}
%         r \rightarrow   &\text{Tasa de rendimiento real}\\
%         \pi \rightarrow &\text{Tasa de Inflacion}\\
%         C_0 \rightarrow &\text{Saldo inical de la cuenta}\\
%         C_t \rightarrow &\text{Saldo Final de la cuenta}\\
%         t \rightarrow &\text{Tiempo}\\
%     \end{align*}
% \end{center}
\section*{Ejemplo ilustrativo}
\textit{Lex Luthor, el hipermagnate inmoviliario, hace 2 meses decidio que deberia revisar su planeacion financiera, ahora mismo necesita revisar cuanto han estado rindiendo sus inversiones bancarias, las cuales le han ofrecido intereses compuestos a un interes del 15\% efectivo, pero como Luthor sabe que este resultado esta contaminado por la inflacion del pais, necesita saber cuanto a sido el rendimiento real de las mismas, para lo cual recupera la informacion de la inflacion de los ultimos 3 años, los datos indican que hubo las siguientes tasas de inflacion anual: 12\%, 6\% y 8\%. ¿Cuanto fue el rendimiento real de las inversiones bancarias de Luthor en cada uno de estos 3 años?}

\subsubsection*{Pasos de la Solucion}
\begin{enumerate}
    \item Volvemos a leer el ejercicio con la intencion de responder la pregunta ¿Que clase de ejercicio es?\\
    \textbf{Es un ejercicio de Inflacion}
    \item Tenemos que trabajar unicamente con tasas de interes efectivas, pero como en este caso ya nos dan efectivas no hay necesidad de realizar ningun procedimiento.
    \item Revisamos la ecuacion de tasas equivalentes $(1+r)=\frac{1+i}{1+\pi}$, en este caso necesitamos  los datos de inflacion e interes para poder obtener nuestros rendimientos reales.
    \item Finalmente ingresamos nuestras diferentes inflaciones, para obtener nuestros rendimientos particulares para cada año.\\
    \textbf{Primer año:}
    \begin{align*}
        (1+r)=\frac{1+15\%}{1+12\%}\\
        r=\frac{1+15\%}{1+12\%}-1\\
        r=\frac{3\%}{1+12\%}\approx2,6786\%
    \end{align*}
    \textbf{Segundo año:}
    \begin{align*}
        (1+r)=\frac{1+15\%}{1+6\%}\\
        r=\frac{1+15\%}{1+6\%}-1\\
        r=\frac{9\%}{1+6\%}\approx8,4906
    \end{align*}
    \textbf{Tercer año:}
    \begin{align*}
        (1+r)=\frac{1+15\%}{1+8\%}\\
        r=\frac{1+15\%}{1+8\%}-1\\
        r=\frac{7\%}{1+8\%}\approx6,4815\%
    \end{align*}
    \text{ }\\
    Con eso realizado el ejercicio se termina.
\end{enumerate}
\subsection{Ejercicios de Capitulo}
\subsection{Problemas para practicar}
\begin{enumerate}
    \item Un inversionista quiere un rendimiento de 12\% por año en un DPF, ¿cual deberia ser la tasa de interes, si se pronostica una inflacion del 4\% durante todo este año?
    \textbf{16,48\%}
    \item Tatiana decidio invertir su dinero en un DPF con un interes del 15\%, si la inflacion en el ultimo año fue del 10\%, ¿cual fue el rendimiento real del DPF?
    \textbf{4,54\%}
    \item María invierte en un bono con una tasa de interés nominal del 9.3\%. Si la inflación durante el período de inversión fue del 3.8\%, ¿cuál fue el rendimiento real del bono?
    \textbf{5,29\%}
    \item Pedro decide depositar su dinero en un fondo mutuo con una tasa de interés nominal del 7.2\%. Si la inflación anual es del 2.5\%, ¿cuál es el rendimiento real del fondo mutuo?    
    \textbf{5,29\%}
    \item Sofía está considerando una cuenta de ahorro, esperando un rendimiento del 6.8\%. Si la tasa de inflación esperada es del 1.2\%, ¿Cual deberia ser la tasa de interes?
    \textbf{8,0816\%}
    \item Luis tiene la opción de invertir esperando un rendimiento del 4\%. Dado que la tasas de interes es del 5.6\%, ¿cuál sería la inflacion?
    \textbf{1,538\%}
    \item Javier decide invertir en un fondo de inversión con una tasa de interés nominal del 10.2\%. Si la inflación esperada es del 3.9\%, ¿cuál sería el rendimiento real del fondo de inversión?
    \textbf{6,064\%}
    \item Marta está considerando un depósito a plazo fijo con una tasa de interés nominal del 7.5\%. Si la inflación proyectada es del 2.3\%, ¿cuál sería el rendimiento real del depósito?
    \textbf{5,083\%}
    \item Diego está evaluando un bono que promete un rendimiento del 9.8\%. Dado que la inflación es del 3.5\%, ¿cuál sería la tasa de interes del bono?
    \textbf{6,087\%}
    \item Carla decide abrir una cuenta de ahorro con una tasa de interés nominal del 5.9\%. Si la inflación durante el período es del 2.8\%, ¿cuál sería el rendimiento real de la cuenta de ahorro?
    \textbf{3,016\%}
    \item Andrés está considerando una inversión que ofrece una tasa de interés nominal del 11.1\%. Si el rendimiento esperado es del 4.9\%, ¿cuál sería la inflacion?
    \textbf{5,91\%}
    \item Laura invierte en bonos con una tasa de interés nominal del 8.9\%. Si la inflación durante el período de inversión es del 3.2\%, ¿cuál sería el rendimiento real de los bonos?
    \textbf{5,523\%}
    \item Carlos decide depositar su dinero en un fondo de inversión con una tasa de interés nominal del 6.5\%. Si la inflación proyectada es del 2.7\%, ¿cuál sería el rendimiento real del fondo de inversión?
    \textbf{3,7\%}
\end{enumerate}
\subsection*{Problemas reto}
\begin{enumerate}
    \item Si Marco Aurelio quiere un rendimiento del 35\% a lo largo de los siguientes 5 años, sabiendo que el pronostico de la inflacion 
    de los siguiente 5 años es del 5\% anual, ¿Cual deberia ser la tasa de interes a la que Marco debe invertir su dinero? ¿Si Marco invierte 
    Bs80.000 cuanto sera la cantidad de retorno? Por variaciones de la inflacion con respecto a las prediciones, Marco Aurelio obtiene un 40\% de rendimiento, 
    si la inflacion fue igual durante los ultimos 5 años, ¿Cual fue la inflacion anual?
    \textbf{56,28\%}\textbf{Bs125.023,5}\textbf{2,224\%}
\end{enumerate}